\section{Phân tích}\label{sec:analysis}

\subsection{Phần bằng chứng còn thiếu nằm ở bước cắt theo ngân sách}\label{sec:residual}
Với RRF(vectơ, BM25), gần như mọi đoạn đáp án đã có trong danh sách ứng viên; phần còn thiếu chủ yếu bị bước cắt loại ra.
Trên 180 câu dev có evidence, \retCandVK\% đoạn đáp án có mặt trong danh sách ứng viên trước khi cắt.
Trong \failVK{} câu không đủ đoạn đáp án sau khi cắt, \cutFailVK{} câu có đoạn đáp án trong ứng viên nhưng bị cắt.
Ở \longAnsCut{}/\cutFailVK{} câu đó, chính đoạn đáp án dài 1.200 token, tức đúng giới hạn độ dài đoạn lúc lập chỉ mục.
Lý do là quy tắc cắt dừng ở đoạn đầu tiên làm tràn ngân sách: vài đoạn hội thoại dài đứng trước có thể chiếm gần hết 4.000 token.
Phát hiện này giải thích vì sao lợi ích biến mất khi hạ ngân sách xuống 2.000 token (\cref{sec:efficiency}), và chỉ ra rằng bước tiếp theo nên là chọn nội dung trong ngân sách chứ không phải tìm thêm ứng viên.

\subsection{Nhóm Null}\label{sec:null}
Mức giảm ở nhóm Null không giải thích được bằng nhãn sai hay cách chấm. Phần lớn nó gắn với việc ngữ cảnh chứa nhiều ``sự thật gần giống'' hơn.
Chúng tôi kiểm ba nguyên nhân.

\paragraph{Nhãn.}
Chúng tôi rà nhãn của 65 câu Null: một câu có đáp án rõ ràng trong kho và hai câu đáng tranh luận.
Loại các câu này chỉ thu hẹp chênh lệch một phần nhỏ \todo{đưa con số chính xác vào script từ \texttt{logs/null\_audit/}}.

\paragraph{Rubric của giám khảo (độ nhạy, không thay số chính thức).}
Một thành viên chấm tay 40 câu Null. Rubric gốc khớp với người ở \humanAgreeOrig{} câu; nhiều câu trả lời ``pha trộn'' (nói không có thông tin rồi vẫn suy đoán) bị chấm \emph{neither}.
Chúng tôi viết một rubric làm rõ cho trường hợp đáp án vàng là ``không đủ thông tin''. Rubric này khớp với người ở \humanAgreeClar{} câu.
Chấm lại 45 câu Null $\times$ 3 lượt sinh $\times$ 3 lượt chấm với rubric mới cho \cref{tab:null}.
Theo rubric làm rõ, các cấu hình không dùng riêng đồ thị gần như không còn cách nhau, nhưng tất cả vẫn có acc thấp hơn và err cao hơn MiniRAG (\cref{tab:null}).
Lời từ chối thuần không bao giờ bị chấm sai (\refusalAcc{} phiếu \emph{accurate}).
Vậy khoảng cách giữa các cấu hình trộn phần lớn nằm ở rubric, còn khoảng cách với MiniRAG là thật.

\begin{table}[t]\centering
\caption{Nhóm Null (45 câu ngoài dev) dưới hai rubric giám khảo, trung bình 3 lượt chấm và các lượt sinh. Rubric làm rõ là phân tích độ nhạy, hiệu chuẩn trên nhãn của một người; số chính thức là cột rubric gốc. Rubric làm rõ gần như không dùng nhãn \emph{neither}, nên chỉ so acc và err giữa hai cột.}
\label{tab:null}
\small
% SINH TỰ ĐỘNG — make_numbers.py
\begin{tabular}{lcc}
\toprule
Cách xếp hạng đoạn & Rubric gốc (chính thức) & Rubric làm rõ (độ nhạy) \\
 & acc / err / neither & acc / err / neither \\
\midrule
MiniRAG (tái lập) & 71{,}1 / 15{,}6 / 13{,}3 & 86{,}7 / 13{,}3 / 0{,}0 \\
RRF(đồ thị, vectơ) & 63{,}0 / 24{,}4 / 12{,}6 & 74{,}8 / 25{,}2 / 0{,}0 \\
Vectơ thuần & 65{,}2 / 27{,}4 / 7{,}4 & 74{,}1 / 25{,}9 / 0{,}0 \\
RRF(đồ thị, vectơ, BM25) & 62{,}2 / 25{,}2 / 12{,}6 & 73{,}3 / 26{,}7 / 0{,}0 \\
RRF(vectơ, BM25) & 57{,}8 / 23{,}0 / 19{,}3 & 74{,}1 / 25{,}9 / 0{,}0 \\
\bottomrule
\end{tabular}

\end{table}

\paragraph{Cơ chế.}
Khi truy xuất tốt hơn, ngữ cảnh chứa nhiều đoạn đúng chủ đề nhưng sai thời điểm hoặc sai sự kiện, và mô hình nhỏ ghép chúng thành một câu trả lời tự tin nhưng sai tiền đề.
Theo phân loại của nhóm, \nullPremise{} câu Null bị chấm \emph{error} thuộc dạng này.
Một tín hiệu ``độ phủ tiền đề'' phân biệt được câu Null với câu có đáp án (AUC \premiseAUC), nhưng dùng nó làm cổng từ chối thì không cứu thêm được câu nào ở mọi ngưỡng.
Lý do: các câu Null có độ phủ thấp thì mô hình vốn đã tự từ chối đúng.
Các hướng sửa khác theo kiểu kiểm chứng hoặc từ chối cũng đều cho kết quả phủ định (\cref{app:tried}).

\subsection{Ví dụ định tính}
\todo{Chọn 1 câu thành công và 1 câu thất bại bằng quy tắc cố định (ví dụ: câu đầu tiên theo thứ tự trong \texttt{nondev435.csv} thuộc mỗi loại), trích câu hỏi, đáp án, đoạn lấy được và câu trả lời của MiniRAG và RRF(vectơ, BM25). Các ví dụ trong \texttt{docs/CAI\_TIEN\_VA\_BENCHMARK.md} thuộc tập dev và đã được chọn tay, không dùng làm ví dụ đại diện.}

\subsection{Độ tin cậy của giám khảo}
Nhiễu giám khảo nhỏ so với nhiễu sinh: độ lệch chuẩn giữa ba lượt chấm là \sdJudgeBase{} điểm, trong khi giữa ba lượt sinh là \sdGenGV{} điểm.
Khoảng \flipRate\% câu đổi phán quyết giữa hai lượt sinh của cùng một cấu hình.
Giám khảo cũng dễ dãi: khi rà tay, \judgeLenient{} câu được chấm \emph{accurate} vẫn chứa ngày, thứ tự hoặc người nói mâu thuẫn với ngữ cảnh.
Sai lệch này tác động lên mọi cấu hình, nhưng chúng tôi chưa đo được nó có tác động như nhau hay không (\cref{sec:limitations}).
